% =====================================================================
%  Chapter 28 review -- the 14-station revision contract.
%  Revision-only: tables, one-liners, pointers back via \secref/\thmref/
%  \defref/\figref/\coderef/\probref.  No re-derivations.
%  Summit: CF 1252F "Regular Forestation" (2400, L5) -- a tree on N <= 4000
%  nodes; a node u (deg >= 2) is a GOOD CUTTING POINT if removing u leaves two
%  or more components that are PAIRWISE IDENTICAL, i.e. isomorphic as unrooted
%  trees; output the maximum number of components obtainable that way, else -1.
%  Pipeline: (i) O(1) component sizes from one root-at-1 pass (Ch.15's
%  complement view n - sz[u]) give the hard filter deg(u) | (N-1) and "every
%  piece has size (N-1)/deg(u)"; (ii) every survivor's pieces are compared as
%  UNROOTED shapes via the chapter's centre form Phi (Ch.12 centre inside the
%  component, one shared interned AHU table for the entire run).  Reuses THREE
%  chapter tools (Phi + the interned labeller + the commutative-hash variant)
%  and TWO earlier-chapter tools (Ch.12 diameter/centre, Ch.15 reroot
%  complement sizes).  Both official samples hand-replayed BEFORE writing:
%  the 13-node tree -> 3 (node 4, three P4 pieces) and the 6-node tree -> -1
%  (N-1 = 5 prime: no degree divides, so zero survivors).  Neither sample
%  separates the wrong-root rung from the correct one -- the two hand-built
%  9-node counterexamples in the thinkbox do, and they are verified here.
%  Exports prob:c22-summit, lst:c22-summit, fig:c22-summit.
%  Mixed set EXACTLY 10, all live-verified this session (titles, ratings and
%  tags re-pulled from the CF problem pages), un-owned, ASCENDING by rating;
%  each hint routes via the shape signature (def:c22-signature) or REJECTS the
%  shape tools and names the true owner (tree DP, DSU-on-tree, path/centroid,
%  HLD/LCT, trie+strings, topological order).
% =====================================================================

\chaptersection{Chapter 28 Summary}

One sentence: \emph{when vertex names, child order and the choice of root carry no
information, the object left behind is a \textbf{shape}, and every question about it --
``are these two trees the same'', ``can the children mirror-pair'', ``how many different
shapes occur'', ``are all the pieces of this cut identical'' -- is one bottom-up fold away:
sort each vertex's child keys, intern the sorted tuple into a fresh id, and compare
integers; unrooted questions add exactly one idea, root at the centre the tree defines for
itself and compare the sorted pair when the centre is an edge
\(\bigl(\secref{c22-ahu}, \thmref{c22-centreiso}, \thmref{c22-classes}\bigr)\).}
Chapter~27 handed this chapter its opening question -- ``which subtree equals which'' --
and this chapter answers it in \(\bigO{n\log n}\) once, for every pair at once.  Each row
below is a pointer, not a proof.

\begin{center}
\renewcommand{\arraystretch}{1.25}
\setlength{\tabcolsep}{3pt}
\begin{tabular}{@{}p{0.20\linewidth} p{0.55\linewidth} p{0.21\linewidth}@{}}
\toprule
\textbf{Result} & \textbf{Statement} & \textbf{Where} \\
\midrule
Interned labelling & Postorder: sort the child ids, hand each fresh sorted tuple a new id;
\((T,r)\cong(T',r')\) iff the root ids match. \(\bigO{n\log n}\), exact, no collisions. &
\thmref{c22-ahu}, \coderef{c22-ahu} \\
Wrong-root trap & The same unrooted tree rooted at a bad vertex gives different ids:
\(R\) rooted at \(1\) is \(4\), its twin \(R'\) rooted at \(1\) is \(7\), yet
\(R\cong R'\). & \figref{c22-twin}, \(\secref{c22-centre}\) \\
Centre form \(\Phi\) & One centre: \(\Phi=\) id at the centre. Centre edge: \(\Phi=\) the
\textbf{sorted pair} of the two endpoint-rooted ids. \(T\cong T'\iff\Phi(T)=\Phi(T')\). &
\defref{c22-canroot}, \thmref{c22-centreiso}, \coderef{c22-unrooted} \\
Mirror symmetry & A rooted tree admits a mirror-pairing child order iff at \emph{every}
vertex at most one child class occurs an odd number of times. & \thmref{c22-sym},
\probref{c22-symmetree} \\
Two hashes & Shape equality needs a \textbf{commutative} combine; path strings need a
\textbf{positional} one, weighted by \(B^{\depth(\cdot)}\) so moving the viewpoint costs
one constant. & \defref{c22-shapehash}, \defref{c22-strhash} \\
Collision ledger & Equal shapes never disagree; different shapes collide with probability
\(\le\depth/2^{64}\). NO is certain, YES is budgeted -- or re-verified by a bijection DFS. &
\thmref{c22-collision} \\
Distinct strings & \(T(v)=\{0\}\cup\bigcup_w\bigl(T(w)\cup\{\mathit{base}_w\}\bigr)\), sets
merged small-to-large: CF 601D's \(f(v)\) with no element ever rewritten. &
\probref{c22-organic}, \coderef{c22-organic} \\
The census & One labelling partitions the vertices: \(|\mathit{tid}|\) distinct shapes,
\(\mathit{freq}[\mathrm{cl}(v)]\) twins. Root-dependent -- \(R\) gives \(4\), the twin
rooted at a leaf gives \(6\). & \thmref{c22-classes}, \figref{c22-classes} \\
The pick & Label-free, order-free, whole-shape object, static tree -- fail one and the
neighbour chapter owns it. & \defref{c22-signature}, \defref{c22-pick} \\
\bottomrule
\end{tabular}
\end{center}

\medskip
{\small\textbf{Algorithms from this chapter.}}
\begin{center}
\renewcommand{\arraystretch}{1.2}
\setlength{\tabcolsep}{2pt}
\begin{tabular}{@{}p{0.33\linewidth} l l >{\raggedright\arraybackslash}p{0.19\linewidth}@{}}
\toprule
\textbf{Algorithm} & \textbf{Prep} & \textbf{Per run} & \textbf{Code} \\
\midrule
Interned AHU labeller (rooted) & \(\bigO{n\log n}\) & \(O(1)\) compare & \coderef{c22-ahu} \\
Unrooted \(\Phi\): centre \(\pm\) centre edge & \(\bigO n\) BFS & \(\bigO{n\log n}\) label & \coderef{c22-unrooted} \\
CF 1800G mirror-symmetry test & \(\bigO{n\log n}\) case & \(O(1)\) per vertex & \coderef{c22-symmetree} \\
Commutative shape hash & \(\bigO{n}\) & \(O(1)\) compare & \coderef{c22-shapehash} \\
CF 601D distinct downward strings & \(\bigO{n\log n}\) exp. & one insert per element & \coderef{c22-organic} \\
Shape census (classes, multiplicities) & \(\bigO{n\log n}\) & \(O(1)\) query & \coderef{c22-classes} \\
Recognition gate / boundary pick & -- & -- & \defref{c22-signature}, \coderef{c22-kit} \\
CF 1252F (Summit): filter then \(\Phi\) & \(\bigO n\) sizes & \(\bigO{n\log n}\) survivor & \coderef{c22-summit} \\
\bottomrule
\end{tabular}
\end{center}

\chaptersection{Master Decision Flowchart}

Post-study ladder: walk top to bottom and stop at the first branch that matches.  It
decides \emph{whether the problem is a shape problem at all} (renaming and reordering
cannot change the answer, the object is a whole downward shape, the tree is static) and,
if so, \emph{which key} to compare: an exact interned id, a commutative hash, a positional
string hash, or the centre-rooted pair.  Everything failing the top gate leaves the
chapter: value payload to tree DP (\(\chref{c10}\)) and DSU-on-tree (\(\chref{c28}\)),
order-bearing traversals to string tools, distances to Ch.~12 (\(\chref{c09}\)) and
centroid (\(\chref{c18}\)), edits to HLD (\(\chref{c16}\)) and link-cut (\(\chref{c23}\)).

\begin{lstlisting}[language={}, basicstyle=\ttfamily\small, numbers=none, frame=single]
  Does RENAMING vertices or SHUFFLING child lists change the answer?
  |
  |-- Yes, values/names matter (sums, colours combined, weights):
  |     +-- NOT this chapter -- tree DP (Ch.13, c10); per-node keyed merge -> DSU
  |         on tree (Ch.17, c28).  Shapes only as ONE component of a joint key.
  |
  |-- Yes, child order matters (left/right roles, preorder strings):
  |     +-- NOT the commutative fold.  Use the traversal string or the positional
  |         hash val(x_i) * B^depth(x_i) (28.2).  CF 1671/E "Preorder" is this shape.
  |
  |-- The object is a PATH or a DISTANCE, not a downward shape:
  |     +-- NOT here.  Ch.12 metrics (c09); pair distributions / path palindromes
  |         -> centroid (Ch.24, c18), where CF 914/E lives.
  |
  |-- Edits (cut/link) arrive between questions:
  |     +-- NOT here.  HLD (Ch.22, c16) or link-cut (Ch.25, c23).  Re-canonicalising
  |         per edit is O(n log n) each -- that is a TLE wearing a correct algorithm.
  |
  |-- STILL HERE: the object is a SHAPE.  Rooted or unrooted?
  |     rooted (a fixed root, subtrees):
  |        compare two trees once      -> interned ids, one SHARED table   (28.1)
  |        thousands of compares       -> commutative hash, O(1) each      (28.2)
  |        pair children into mirrors  -> one odd class per vertex         (28.3, 1800G)
  |        how many shapes / class size-> |tid| and freq[cl[v]]            (28.4)
  |        distinct downward strings   -> positional hash + small-to-large (601D)
  |     unrooted (names and root are noise):
  |        root at the CENTRE (Ch.12).  One centre -> one id.  Centre EDGE ->
  |        compare the SORTED PAIR of the two endpoint-rooted ids; picking one
  |        endpoint is the false-NO trap (28.3, fig:c22-pair).
  |        every component of one cut  -> size filter, then Phi per piece  (Summit)
  |
  |-- Judge prints a decisive YES/NO and one collision would cost you the round:
  |     +-- take the exact route (ids), or re-verify every hash-YES with an explicit
  |         bijection DFS (28.2).  Never pay the collision budget twice by accident.
\end{lstlisting}

\chaptersection{Common Mistakes Checklist}

\begin{itemize}[nosep, leftmargin=1.4em]
  \item \textbf{Rooting an unrooted question at vertex \(1\).} \(R\) and its twin \(R'\) are
        the same eight edges, yet root-at-\(1\) gives \(4\) against \(7\)
        (\(\figref{c22-twin}\)).  Root at the centre (\(\thmref{c22-centreiso}\)) -- any
        rule that mentions a label can be broken by renaming.
  \item \textbf{Choosing one endpoint of the centre edge.} With a centre edge the tree
        cannot tell the endpoints apart, so a single rooted form is not a function of the
        tree: \(\Phi(Q_1)=\langle C_2,\langle L,C_2\rangle\rangle\) against
        \(\Phi(Q_3)=\langle L,C_2,C_3\rangle\) at the same endpoint, while the trees are
        isomorphic (\(\figref{c22-pair}\)).  Compare the \emph{sorted pair}.
  \item \textbf{A fresh \(\mathit{tid}\) table per tree.} Ids are only comparable inside one
        table: intern both trees (and both rootings of the centre-edge case) against
        \textbf{one} shared map, or the equality of two integers means nothing
        (\(\secref{c22-ahu}\), pitfall~1).
  \item \textbf{Pairing children by subtree \emph{size} in 1800G.} A three-node chain and a
        three-node star have equal sizes and different shapes; the pairing key must be the
        shape id or a shape hash (\(\thmref{c22-sym}\)).
  \item \textbf{Testing symmetry only at the root.} The odd-class condition must hold at
        \emph{every} vertex -- official samples 4 and 5 of CF 1800G fail exactly one level
        down (\(\probref{c22-symmetree}\)).
  \item \textbf{Mixing the two hashes.} The commutative fold identifies
        \(\langle A,B\rangle\) with \(\langle B,A\rangle\); the positional one does not.
        Using the first on strings collapses distinct preorder readings into one collision
        you cannot lose (\(\defref{c22-shapehash}\), \(\defref{c22-strhash}\)).
  \item \textbf{Spending a probabilistic YES on a decisive answer.} A single \(64\)-bit
        hash is wrong with probability \(\le\depth/2^{64}\) per compared pair -- fine for
        a batch, not fine when one YES costs the round.  Double, or verify by an explicit
        bijection DFS (\(\thmref{c22-collision}\)).
  \item \textbf{Counting vertices instead of classes.} The census wants
        \(|\mathit{tid}|\), not \(n\); and multiplicities are only meaningful inside one
        table (\(\thmref{c22-classes}\), \(\secref{c22-distinct}\)'s warnings).
  \item \textbf{Recursion overflow and \(32\)-bit overflow in 601D.} \(n=3\cdot10^5\) kills
        the call stack; \(f(v)+c_v\) needs \code{long long}.  Both are WA/RE with a correct
        algorithm (\(\coderef{c22-organic}\)).
  \item \textbf{Reusing stale ids after an edit.} One cut changes the centre, the postorder
        and every id above it; canonical forms are valid only until the topology changes
        (\(\defref{c22-signature}\), test~4).
\end{itemize}

\chaptersection{WA Clinic}

Verdicts are symptoms; the causes are the short list this chapter drilled.  Diagnose
top-down -- every leaf names a listing.

\begin{figure}[H]
\centering
\begin{tikzpicture}[diagram,
    v/.style={boxnode, fill=primaryblue!12, draw=primaryblue, text width=22mm, align=center},
    b/.style={boxnode, fill=difforange!12, draw=difforange, text width=12mm},
    l/.style={boxnode, fill=successgreen!10, draw=successgreen, text width=48mm,
              font=\scriptsize, align=left}]
  \node[v] (verdict) {\textbf{Verdict on the judge}};
  \node[b, below left=8mm and 5mm of verdict] (wa) {\textbf{WA}};
  \node[b, below right=8mm and 5mm of verdict] (tle) {\textbf{TLE / RE}};
  \node[l, below right=1mm and 2mm of wa.south] (wa1)
    {unrooted question rooted at \(1\), or one endpoint of the centre edge picked --
     \(R\) against \(R'\) says NO, \((Q_1,Q_3)\) says NO (\(\figref{c22-twin}\),
     \(\figref{c22-pair}\))};
  \node[l, below left=1mm and 2mm of wa.south] (wa2)
    {1800G paired by subtree size, or symmetry checked at the root only; 601D sets keyed by
     the commutative hash instead of the positional one (\(\coderef{c22-symmetree}\),
     \(\coderef{c22-organic}\))};
  \node[l, below right=1mm and 2mm of tle.south] (tle1)
    {AHU re-run per compare (\(\bigO{n}\) each) instead of one hash precompute, or the
     601D merge copying the big child instead of adopting it (\(\defref{c22-shapehash}\),
     \(\secref{c22-hash}\))};
  \node[l, below left=1mm and 2mm of tle.south] (tle2)
    {recursive labeller or recursive rootify at \(n=3\cdot10^{5}\) -- RE; go iterative, and
     keep \(f+c\) in \code{long long} (\(\coderef{c22-organic}\), \(\coderef{c22-unrooted}\))};
  \draw[flow] (verdict) -- (wa);
  \draw[flow] (verdict) -- (tle);
  \draw[flow] (wa) -- (wa1);
  \draw[flow] (wa) -- (wa2);
  \draw[flow] (tle) -- (tle1);
  \draw[flow] (tle) -- (tle2);
\end{tikzpicture}
\caption*{WA: the root or the centre-edge endpoint was chosen by a rule the tree does not
define, or the pairing key was size instead of shape. TLE or RE: re-canonicalising per
compare, a merge that copies, and recursion on a \(3\cdot10^{5}\)-vertex tree.}
\end{figure}

\chaptersection{Quick Reference Card}

\begin{quickref}[Chapter 28 at a glance]
\textbf{Tree isomorphism \(=\) remove every accidental freedom (names, child order, root
choice) and then compare integers.  Interned ids are exact and cost \(\bigO{n\log n}\) once;
commutative hashes cost \(\bigO n\) and \(O(1)\) per compare but charge a collision
premium; the unrooted case is rooted at the centre and, for a centre edge, compared as a
sorted pair.}

\medskip
\textbf{Rooted AHU (28.1).} Postorder; \(t\gets\) sorted child ids; \(\mathit{tid}[t]\) or a
fresh id.  \((T,r)\cong(T',r')\iff\mathrm{cl}(r)=\mathrm{cl}(r')\), one shared table
(\(\thmref{c22-ahu}\), \(\coderef{c22-ahu}\)).  On \(R\): leaf \(1\), star \(2\), chain
\(3\), root \(4\).

\medskip
\textbf{Hashes (28.2).} Shape: \(h(u)=1+\sum_w\mathrm{rho}(\depth(w),h(w))\) -- commutative,
so child order cannot matter (\(\defref{c22-shapehash}\)).  Strings:
\(G=\sum_i\mathit{val}(x_i)B^{\depth(x_i)}\), start letter excluded, so a viewpoint shift is
one constant (\(\defref{c22-strhash}\)).  Ledger: NO certain, YES wrong
\(\le\depth/2^{64}\) (\(\thmref{c22-collision}\)).

\medskip
\textbf{Unrooted and mirror (28.3).} Centre from the diameter (Ch.12):
\(\Phi\) is one id, or the sorted pair of two.  Symmetry: at most one odd child class per
vertex, keyed by shape id (\(\thmref{c22-centreiso}\), \(\thmref{c22-sym}\),
\(\coderef{c22-unrooted}\), \(\coderef{c22-symmetree}\)); CF 1800G's six samples
\(\to\) YES NO YES NO NO YES.

\medskip
\textbf{Counting (28.4, 601D).} \(|\mathit{tid}|\) distinct shapes,
\(\mathit{freq}[\mathrm{cl}(v)]\) twins; per-node distinct-string sets merged
small-to-large, \(T(v)=\{0\}\cup\bigcup_w(T(w)\cup\{\mathit{base}_w\})\)
(\(\thmref{c22-classes}\), \(\probref{c22-organic}\)).

\medskip
\textbf{The pick (28.5).} Label-free \emph{and} order-free \emph{and} whole-shape \emph{and}
static \(\Rightarrow\) this chapter; values \(\to\) \chref{c10}/\chref{c28}; order \(\to\)
traversal strings; distances \(\to\) \chref{c09}/\chref{c18}; edits \(\to\)
\chref{c16}/\chref{c23} (\(\defref{c22-signature}\), \(\coderef{c22-kit}\)).
\end{quickref}

\chaptersection{Write-From-Memory Drills}

Five signature skeletons, each lifted from a listing already in the chapter.  Cover the
marked line, reproduce from memory, then check the answer one line below.

\begin{lstlisting}[caption={Drill 1 -- the interned AHU labeller (from lst:c22-ahu)}]
int label(int u){ vector<int> t;
    for(int v : kids[u]) t.push_back(label(v));
    sort(t.begin(), t.end());                       // order is noise: kill it
    map<vector<int>,int>::iterator it = tid.find(t);
    if(it != tid.end()) return cl[u] = it->second;  // recall an existing shape
    int id = nextId++; tid[t] = id; return cl[u] = id;   // fresh shape name
}
\end{lstlisting}
\noindent{\scriptsize Key: postorder so every child is named before its parent; \texttt{tid}
is shared by both trees \emph{and} both rootings; equal id \(\iff\) isomorphic, and the sort
is what makes child order irrelevant (\(\thmref{c22-ahu}\)).}

\begin{lstlisting}[caption={Drill 2 -- centre form Phi of an unrooted tree (from lst:c22-unrooted)}]
A = farthest(start, d);  B = farthest(A, d);  len = d[B];   // diameter (Ch.12)
rootAt(A);  m = B;  for(i = 0; i < len/2; i++) m = par[m];  // walk to the middle
if(len % 2 == 0){ rootAt(m); return {cl[m], cl[m]}; }       // ONE centre
c1 = m;  c2 = par[m];                                       // centre EDGE
rootAt(c1); x = label(c1);  rootAt(c2); y = label(c2);      // same tid table
return {min(x,y), max(x,y)};                                // compare the PAIR
\end{lstlisting}
\noindent{\scriptsize Key: two centre endpoints are indistinguishable, so keep both rooted
forms and sort them; picking one endpoint loses real isomorphisms (\(\figref{c22-pair}\),
\(\thmref{c22-centreiso}\)).}

\begin{lstlisting}[caption={Drill 3 -- the mirror-symmetry test (from lst:c22-symmetree)}]
sort(t.begin(), t.end());                       // t = child shape ids of u
int odd = 0;
for(j = 0; j < t.size(); ){ k = j; while(k < t.size() && t[k] == t[j]) k++;
    if((k - j) & 1) odd++;  j = k; }            // one run of odd length = odd class
if(odd > 1) ok = false;                         // at EVERY vertex
\end{lstlisting}
\noindent{\scriptsize Key: pairs are mirrors, so at most one class may be left unpaired --
and the middle child must itself be symmetrical, which is why the test is per vertex, not per
root (\(\thmref{c22-sym}\), CF 1800G).}

\begin{lstlisting}[caption={Drill 4 -- distinct downward strings, 601D core (from lst:c22-organic)}]
base[u] = (ull)(s[u]-'a'+1) * pw[dep[u]];       // depth-weighted: shift = +one term
for(u in reverse postorder){  big = child with the largest bag;
    bag[u] = (big==-1) ? new unordered_set<ull>() : bag[big];   // ADOPT, do not copy
    bag[u]->insert(0ULL);
    for(v in kids[u]) if(v != big){ bag[u]->insert(base[v]);
        for(e : *bag[v]) bag[u]->insert(e);  delete bag[v]; }   // pour the small ones
    f[u] = bag[u]->size(); }
\end{lstlisting}
\noindent{\scriptsize Key: no element is ever rewritten, so a set survives being adopted;
each element moves \(O(\log n)\) times -- \(\bigO{n\log n}\) expected
(\(\secref{c22-hash}\), \(\probref{c22-organic}\)).}

\begin{lstlisting}[caption={Drill 5 -- the Summit's filter and per-piece form (from lst:c22-summit)}]
d = adj[u].size();  if(d < 2 || (n-1) % d) continue;          // sizes must divide
s = (n-1) / d;
for(v in adj[u]) if(sizeOfPiece(u,v) != s) goto reject;       // O(1) each, from sz[]
for(v in adj[u]) if(phiOfPiece(u,v) != first) reject;         // centre + shared tid
best = max(best, d);
// sizeOfPiece: v child -> sz[v];  v = par[u] -> n - sz[u]   (Ch.15 complement view)
\end{lstlisting}
\noindent{\scriptsize Key: the cheap invariant (equal component \emph{size}) is a filter,
never a proof; equality of unrooted shapes is \(\Phi\), computed per piece with the
\emph{same} table all run (\(\probref{c22-summit}\)).}

\chaptersection{Complete Problem Checklist}

Every problem this chapter offers -- two fully solved in-section, worked traces on the
running example \(R\) and its twins, and one Summit on its own CF samples.  Tick \emph{only}
after an unaided solve; a hinted attempt counts as half.

\begin{center}
\footnotesize
\setlength{\tabcolsep}{4pt}
\renewcommand{\arraystretch}{1.25}
\begin{tabular}{@{}c p{0.16\linewidth} l p{0.20\linewidth} >{\raggedright\arraybackslash}p{0.26\linewidth} c@{}}
\toprule
\textbf{\#} & \textbf{Problem} & \textbf{Source} & \textbf{Key technique} & \textbf{What to look for} & \textbf{Done?} \\
\midrule
\multicolumn{6}{l}{\textbf{Chapter opener -- what is a shape?}} \\
W28.0 & the look-alike pair \((T_1,T_2)\) & worked & \lv{3} invariants vs.\ distance profile & same size, degree sequence, diameter, centre edge -- still NO (\(\figref{c22-trap}\)) & \(\square\) \\
\midrule
\multicolumn{6}{l}{\textbf{Section 28.1 -- rooted AHU}} \\
W28.1 & label \(R\), then its twin \(R'\) & worked & \lv{4} sort + intern, one shared table & \(R\): leaf \(1\), star \(2\), chain \(3\), root \(4\); \(R'\) rooted at \(1\) reaches \(7\) & \(\square\) \\
\midrule
\multicolumn{6}{l}{\textbf{Section 28.2 -- hashing}} \\
W28.2a & commutative vs.\ positional hash & worked & \lv{4} two combines, one ledger & shape fold ignores order; \(B^{\depth}\) strings do not; YES costs \(2^{-64}\) & \(\square\) \\
W28.2b & Acyclic Organic Compounds & CF 601D & \lv{5} depth-weighted sets + small-to-large & \(T(v)=\{0\}\cup\bigcup_w(T(w)\cup\{\mathit{base}_w\})\); samples \(\to\) 51/3 and 6/2 & \(\square\) \\
\midrule
\multicolumn{6}{l}{\textbf{Section 28.3 -- centre and mirrors}} \\
W28.3a & \(\Phi\) on \((R,R')\), \((T_1,T_2)\), \(P_6\), \((Q_1,Q_3)\) & worked & \lv{4} centre-edge sorted pair & \(4{=}4\) YES; \((5,6)\neq(9,10)\) NO; \(P_6\) self-dual; single endpoint lies & \(\square\) \\
W28.3b & Symmetree & CF 1800G & \lv{4} one odd child class per vertex & key by shape not size; six samples \(\to\) YES NO YES NO NO YES & \(\square\) \\
\midrule
\multicolumn{6}{l}{\textbf{Section 28.4 -- counting shapes}} \\
W28.4 & shape census of \(R\) and \(R'\) & worked & \lv{4} \(|\mathit{tid}|\) and \(\mathit{freq}[\mathrm{cl}]\) & \(4\) classes (sizes \(5,2,1,1\)) against \(6\) -- the census is root-dependent & \(\square\) \\
\midrule
\multicolumn{6}{l}{\textbf{Section 28.5 -- recognition / boundary (worked only)}} \\
W28.5 & four specs, four verdicts & worked & \lv{5} the four tests & census keeps it; values, order (CF 1671/E) and edits each hand it on & \(\square\) \\
\midrule
\multicolumn{6}{l}{\textbf{Summit (chapter end)}} \\
W28 & Regular Forestation & CF 1252F & \lv{5} size filter, then \(\Phi\) per piece & \(\deg(u)\mid N{-}1\); pieces compared unrooted; samples \(\to\) 3 and \(-1\) & \(\square\) \\
\bottomrule
\end{tabular}
\end{center}

\chaptersection{Recommended Study Order}

Four focused days. Day~1 the labeller and the hand run on \(R\); day~2 the two hashes and
CF~601D; day~3 the centre form and CF~1800G; day~4 the census, the pick, the Summit and the
first mixed entries.

\begin{center}
\small
\renewcommand{\arraystretch}{1.25}
\setlength{\tabcolsep}{3pt}
\begin{tabular}{@{}llp{0.55\linewidth}@{}}
\toprule
\textbf{Day} & \textbf{Focus} & \textbf{Tasks} \\
\midrule
1 & Read + do 28.1
  & label \(R\) bottom-up by hand to leaf \(1\), star \(2\), chain \(3\), root \(4\)
    (\(\coderef{c22-ahu}\)); then label the twin \(R'\) rooted at \(1\) and watch it reach
    \(7\) (\(\figref{c22-twin}\)); say out loud why both trees must be interned in
    \emph{one} \(\mathit{tid}\) table before any equality is meaningful
    (\(\thmref{c22-ahu}\), \(\defref{c22-canlab}\)) \\
2 & Read + do 28.2
  & write both combines from memory and name which one collapses
    \(\langle A,B\rangle\) with \(\langle B,A\rangle\) (\(\defref{c22-shapehash}\),
    \(\defref{c22-strhash}\)); state the ledger -- NO certain, YES
    \(\le\depth/2^{64}\) (\(\thmref{c22-collision}\)); then CF~601D, replaying the adopt
    and pour to \(51/3\) and \(6/2\) (\(\probref{c22-organic}\)) \\
3 & Read + do 28.3
  & compute \(\Phi\) on \((R,R')\), on \(P_6\) and on the false pair \((T_1,T_2)\), then on
    \((Q_1,Q_3)\) to feel the centre-edge trap (\(\defref{c22-canroot}\),
    \(\thmref{c22-centreiso}\), \(\coderef{c22-unrooted}\)); then CF~1800G, six samples to
    YES NO YES NO NO YES (\(\thmref{c22-sym}\), \(\probref{c22-symmetree}\)) \\
4 & Read + do 28.4, 28.5, Summit
  & census \(R\) and \(R'\): \(4\) classes against \(6\), sizes \(5,2,1,1\)
    (\(\thmref{c22-classes}\)); run the four specs through the pick
    (\(\defref{c22-signature}\), \(\coderef{c22-kit}\)); then CF~1252F, both samples
    (\(\probref{c22-summit}\)); close with the first three mixed entries \\
\bottomrule
\end{tabular}
\end{center}

\chaptersection{Difficulty Progression}

Tiers are competency gates -- each row names the \emph{skill} the tier unlocks, not a
rating number.

\begin{center}
\small
\renewcommand{\arraystretch}{1.3}
\setlength{\tabcolsep}{3pt}
\begin{tabular}{@{}lp{0.30\linewidth}p{0.42\linewidth}@{}}
\toprule
\textbf{Tier} & \textbf{Problems} & \textbf{What passing this tier proves} \\
\midrule
\lv{3} \ \(\bigstar\bigstar\bigstar\)
  & the hand labelling of \(R\) and \(R'\); \(\Phi(P_6)\); the look-alike pair
    \((T_1,T_2)\)
  & you can produce canonical names without a compiler, and you can list four invariants
    that are \emph{not} a complete certificate (\(\figref{c22-trap}\)) \\
\lv{4} \ \(\bigstar\bigstar\bigstar\bigstar\)
  & CF 1800G, the AHU listing, the unrooted compare
  & you code the exact route -- shared table, centre root, sorted pair for the centre edge
    -- and you test symmetry at every vertex, keyed by shape (\(\coderef{c22-ahu}\),
    \(\coderef{c22-symmetree}\)) \\
\lv{5} \ \(\bigstar\bigstar\bigstar\bigstar\bigstar\)
  & CF 601D, the shape census, the boundary pick
  & you combine a hash with a merge discipline (nothing is ever rewritten), and you
    \emph{reject} the shape tools when a value, an order or an edit enters the spec
    (\(\defref{c22-pick}\)) \\
\lv{5} \ \(\bigstar\bigstar\bigstar\bigstar\bigstar\)
  & CF~1252F (the Summit), mixed at CF~2300--2400
  & you use a cheap invariant as a \emph{filter} and an exact canonical form as the
    \emph{proof}, per piece, with one table for the whole run
    (\(\probref{c22-summit}\)) \\
\bottomrule
\end{tabular}
\end{center}

\chaptersection{Reduction Gallery}

Every entry is ``one question shape \(\to\) one move (or an honest hand-off).'' Each line
collapses a problem family onto a listing already in the chapter.

\begin{itemize}[nosep, leftmargin=1.4em]
  \item ``are these two \emph{rooted} trees the same tree?'' \(\to\) postorder, sort child
        ids, intern \((T,r)\cong(T',r')\iff\mathrm{cl}(r)=\mathrm{cl}(r')\), one shared
        table (\(\thmref{c22-ahu}\), \(\coderef{c22-ahu}\)).
  \item ``which subtree of \(T\) equals which subtree of \(T'\)'' -- Ch.~27's parting
        question \(\to\) one labelling of both trees answers every pair at once, \(O(1)\)
        per pair (\(\defref{c22-canlab}\); the baton in \(\secref{c22-ahu}\)).
  \item ``are these two \emph{unrooted} trees the same?'' \(\to\) root at the centre
        (Ch.12, black box), compare one id, or the \textbf{sorted pair} for a centre edge
        (\(\thmref{c22-centreiso}\), \(\coderef{c22-unrooted}\)).
  \item ``can the children be ordered so the tree mirrors itself?'' \(\to\) at most one
        child class of odd multiplicity at \emph{every} vertex, keyed by shape id
        (\(\thmref{c22-sym}\), \(\probref{c22-symmetree}\)); CF~1182/D runs the same test
        with values attached (owned by Ch.~20, prose cite only).
  \item ``how many different shapes occur / how many twins does this one have?'' \(\to\)
        \(|\mathit{tid}|\) and \(\mathit{freq}[\mathrm{cl}(v)]\) from the same pass
        (\(\thmref{c22-classes}\), \(\coderef{c22-classes}\)).
  \item ``\(10^5\) equality queries, static tree'' \(\to\) commutative shape hash,
        \(\bigO n\) precompute, \(O(1)\) compare, and the ledger decides how much YES you
        are buying (\(\defref{c22-shapehash}\), \(\thmref{c22-collision}\)).
  \item ``distinct downward strings from every vertex'' (CF~601D) \(\to\) depth-weighted
        positional hash, sets merged small-to-large so no element is rewritten
        (\(\probref{c22-organic}\), \(\coderef{c22-organic}\); the union itself is Ch.~17,
        \(\chref{c28}\)).
  \item ``count distinct preorder readings after child swaps'' (CF~1671/E) \(\to\) the
        \emph{order-bearing} string hash, plus the test ``do these two readings concatenate
        to the same string'' -- equality of strings, never of shapes
        (\(\defref{c22-strhash}\); \(\secref{c22-boundary}\)'s order-free caveat).
  \item ``every piece of one cut identical'' (CF~1252F) \(\to\) \(\deg(u)\mid(N{-}1)\) and
        equal piece sizes as a filter, then \(\Phi\) per piece as the proof
        (\(\probref{c22-summit}\), \(\coderef{c22-summit}\)).
  \item ``the shape is the same but the answer combines \emph{values}'' \(\to\)
        \emph{not} here: bottom-up numeric folds are tree DP (\(\chref{c10}\)); a per-node
        keyed answer is DSU-on-tree (\(\chref{c28}\)). Shape may appear only as one half of
        a joint key (\(\defref{c22-signature}\), test~1).
  \item ``the object is a path, a distance, or a letter string along one'' \(\to\)
        \emph{not} here: Ch.~12 metrics (\(\chref{c09}\)); pair and path distributions with
        no query list are centroid's, which is where CF~914/E lives (\(\chref{c18}\)).
  \item ``cuts and links arrive between questions'' \(\to\) \emph{not} here: HLD
        (\(\chref{c16}\)) or link-cut (\(\chref{c23}\)); re-canonicalising the form after
        every single edit is a \(\bigO{n\log n}\) tax on an otherwise correct algorithm
        (\(\defref{c22-signature}\), test~4).
  \item ``the input is a general graph'' \(\to\) the centre trick has no anchor and
        interning loses information: colour refinement, or \(k\)-WL for larger \(k\), is the
        honest generalisation -- strictly weaker than tree isomorphism
        (\(\defref{c22-pick}\)).
\end{itemize}

\chaptersection{Mixed Problem Set}

Ten entries, each live-verified against the current Codeforces problem set (titles,
ratings and tags re-pulled from the \texttt{problemset.problems} API and the problem pages
this session) and absent from every other roster in the book. None says ``test
isomorphism''; half of them deliberately \emph{do not} want the shape tools, and the drill
is to run each through the four tests (\(\defref{c22-signature}\)) and say so. Each hint
names the reduction -- or the rejection and the chapter that owns it -- ordered ascending
by rating.

\begin{mixedset}
  \practice{DFS Checker (Easy Version)}{CF 2002D1}
    {https://codeforces.com/problemset/problem/2002/D1}{CF 1900 L4}
    {Tags \emph{brute force, data structures, dfs, graphs, hashing, trees}. A perfect
     binary tree on \(n=2^k-1\) vertices, a permutation \(p\) of its labels, and \(q\)
     persistent swaps \(p_x\!\leftrightarrow\!p_y\); after each, decide whether \(p\) is
     \emph{some} DFS order of the tree. Test~2 of the signature fires backwards here:
     shuffling children is exactly what \emph{defines} the answer set, so an order-free
     fold cannot be the key -- \(p\) is a DFS order iff every subtree occupies one
     \emph{contiguous} block of positions \emph{starting} at its own vertex, a laminarity
     statement about intervals (\(\chref{c24}\)'s \(\tin/\tout\) is the static version of
     the same fact). Reject \(\Phi\); keep the interval view and brute force it in the easy
     version. The preorder-string family it belongs to is CF~1671/E below.}
  \practice{Autocompletion}{CF 928D}
    {https://codeforces.com/problemset/problem/928/D}{CF 1900 L4}
    {Tags \emph{strings, trees}. Arcady types a text; whenever the on-screen prefix of the
     current word matches exactly one already-typed word, the editor offers it and one
     click jumps to it; minimise the clicks. Test~1 fails immediately: the children of a
     node are \emph{letters}, so renaming them changes the answer -- the tree here is a
     trie, and ordered-labelled structure is the whole point. The DP runs along the trie in
     first-appearance order (prefix uniqueness is a subtree-of-trie count). Owner: string
     / trie DP; \(\defref{c22-signature}\) sends it away, and \(\defref{c22-strhash}\) is
     only the machinery you would use to name trie nodes cheaply.}
  \practice{Weight the Tree}{CF 1646D}
    {https://codeforces.com/problemset/problem/1646/D}{CF 2000 L4}
    {Tags \emph{constructive algorithms, dfs, dp, implementation, trees}. Assign positive
     integer weights; a vertex is \emph{good} when its weight equals the sum of its
     neighbours' weights; maximise the number of good vertices, then minimise the total
     sum, and print a witness. The objective folds \emph{numbers}, so test~1 fails -- the
     answer for a vertex is not a function of its shape alone. The shape instinct is not
     useless (identical branches admit identical optimal patterns, which is why a
     constant-size state DP per vertex suffices), but the decision is a bottom-up numeric
     fold: tree DP (\(\chref{c10}\)), with \(\Phi\) only as a way to sanity-check that your
     states are shape-invariant.}
  \practice{Preorder}{CF 1671E}
    {https://codeforces.com/problemset/problem/1671/E}{CF 2100 L5}
    {Tags \emph{combinatorics, divide and conquer, dp, dsu, hashing, sortings, trees}. A
     perfect binary tree on \(2^n-1\) vertices (\(n\le 18\)) labelled A/B; you may swap the
     two children of any non-leaf any number of times; count the \emph{distinct} preorder
     strings obtainable, mod \(998244353\). Official samples: \texttt{BAAAAAAAABBABAB}
     (\(n{=}4\)) \(\to\) \(16\); \texttt{BAA} \(\to\) \(1\); \texttt{ABA} \(\to\) \(2\);
     \texttt{AAB} \(\to\) \(2\); \texttt{AAA} \(\to\) \(1\). This is the chapter's
     order-bearing twin: the object is a \emph{string}, so the positional hash
     \(\mathit{val}(x_i)B^{\depth(x_i)}\) is the right key (\(\defref{c22-strhash}\)) and
     the commutative fold would collapse the two orders into one and lose the count
     (\(\secref{c22-hash}\)'s mixing warning). The double-count correction at a node is the
     test ``do reading \(a\) of the left half and reading \(b\) of the right half satisfy
     \(ab=ba\)?'' -- \(\texttt{BAA}\to1\) is exactly that case. Budget the collision
     ledger (\(\thmref{c22-collision}\)) against \emph{under}-counting.}
  \practice{Hossam and (sub-)palindromic tree}{CF 1771D}
    {https://codeforces.com/problemset/problem/1771/D}{CF 2100 L5}
    {Tags \emph{brute force, data structures, dfs, dp, strings, trees}. Letters in
     vertices; \(s(v,u)\) is the string along the unique path \(v\to u\); a
     sub-palindrome is a palindromic subsequence; report the length of the longest maximal
     sub-palindrome over \emph{all} paths (\(n\le 2000\)). Two tests fail at once: the
     object is a path, not a downward shape, and reversal order matters. Shape hashing
     cannot see a subsequence. Owner: path reasoning with letter counts -- Ch.~12 metrics
     (\(\chref{c09}\)) for the path machinery, and the ``split at the centre of a path''
     discipline is centroid's (\(\chref{c18}\)), which is where CF~914/E lives. Nearest
     shape-flavoured fact: the answer depends on the multiset of letters on the path, not
     their order -- a \emph{commutative} fold, but over letters, not over shapes.}
  \practice{Valid Sets}{CF 486D}
    {https://codeforces.com/problemset/problem/486/D}{CF 2100 L5}
    {Tags \emph{dfs, dp, math, trees}. Given \(d\) and vertex values \(a_i\), count the
     non-empty connected sets \(S\) whose values satisfy \(\max_S a-\min_S a\le d\), mod
     \(10^9{+}7\) (\(n,d\le 2000\)). Test~1 fails flat -- values decide membership, and the
     statement's own note says \(\{1,4\}\) passes the value test but fails connectedness.
     The standard reduction is to root the tree at the set's minimum (tie-broken by index,
     which is the classic double-count repair) and fold downward inside the value window:
     tree DP (\(\chref{c10}\)). Nothing here is invariant under renaming, so
     \(\defref{c22-signature}\) rejects it at the first test; use the rejection as a
     calibration, not a loss.}
  \practice{Ants in Leaves}{CF 622E}
    {https://codeforces.com/problemset/problem/622/E}{CF 2200 L5}
    {Tags \emph{dfs, greedy, sortings, trees}. One ant per leaf; each second any subset of
     ants moves to their parent; no two ants may occupy a vertex at once except the root;
     minimise the time to gather all of them (\(n\le 5\cdot10^5\)). The payload is
     \emph{depth}, so test~1 fails; the answer is a greedy pairing of ant arrival times
     inside each child list, sorted, with the root's children handled jointly because the
     root tolerates collisions. Compare with \(\thmref{c22-sym}\): both ask you to
     \emph{pair} things per vertex, and the pairing key there is a shape id while here it is
     a number -- a clean illustration of why the shape signature is about \emph{which}
     quantities are invariant, not about how tidy the local rule looks.}
  \practice{DFS Checker (Hard Version)}{CF 2002D2}
    {https://codeforces.com/problemset/problem/2002/D2}{CF 2300 L5}
    {Tags \emph{binary search, data structures, dfs, graphs, hashing, trees}. The same
     decision for a \emph{generic} rooted tree, \(n\le 3\cdot10^5\), \(q\le 10^5\)
     persistent swaps. Now the laminarity view has to survive edits: maintain, per vertex,
     the span of positions its subtree occupies and the number of ``broken'' vertices, so
     the answer is a single counter equal to zero -- a set/Fenwick job on
     \(\tin/\tout\)-shaped information (\(\chref{c24}\), \(\chref{c14}\)). The
     \emph{hashing} tag is honest but easy to misread: what you hash is a \emph{set of
     vertices} (a commutative fold over labels, order-free by construction), not a shape;
     \(\defref{c22-shapehash}\) supplies the mechanism, while \(\defref{c22-signature}\)
     still rejects the problem -- child order is the answer set, and the swaps make the
     reading dynamic.}
  \practice{Trees of Tranquillity}{CF 1528C}
    {https://codeforces.com/problemset/problem/1528/C}{CF 2300 L5}
    {Tags \emph{data structures, dfs, greedy, trees}. Two labelled rooted trees on the same
     vertex set; the memorial graph joins \(u,v\) iff one is an ancestor of the other in
     Soroush's tree and \emph{neither} is an ancestor of the other in Keshi's; find the
     maximum clique. Test~1 fails in the harshest way available: the two trees share the
     alphabet of labels, so renaming anything destroys the relation between them -- this is
     a question about a \emph{pair of partial orders}, not about shape. Reduce to DFS on
     \(T_1\) carrying an ordered set of \(\tin\) values from \(T_2\), greedy replacement
     (Ch.~5's tour plus a set, \(\chref{c24}\)). The lesson to state out loud: identical
     \emph{shapes} for both trees would tell you nothing.}
  \practice{Plan of Lectures}{CF 1463E}
    {https://codeforces.com/problemset/problem/1463/E}{CF 2400 L5}
    {Tags \emph{constructive algorithms, dfs, dsu, graphs, implementation, sortings, trees}.
     \(n\) topics, a prerequisite map \(p_i\) forming a rooted forest (exactly one topic
     has none), and \(k\) special pairs \((x_i,y_i)\) that must appear as
     \emph{consecutive} lectures; decide feasibility and print an order. Nothing is
     invariant under renaming -- topics carry identities and the constraint mentions them
     -- so tests~1 and~2 both fail, and the object is a linear extension with adjacency
     obligations, not a shape. Reduce: the forced pairs must form vertex-disjoint directed
     paths consistent with \(p\), contract each into a chain (DSU / degree check), then DFS
     the contracted forest for the order. Owner: topological construction plus DSU
     (\(\chref{c28}\)'s family), \emph{not} this chapter; take it as the set's cleanest
     rejection at maximum rating.}
\end{mixedset}

\chaptersection{The Summit -- CF 1252F Regular Forestation}

\begin{notebox}[Why this is the chapter top rope]
Everything before it compared \emph{two} trees. The Summit compares \(d\) trees at once,
\emph{inside} one tree, at a location the input does not name -- and it wants the best
location. Three separate difficulties collapse onto three tools of this chapter, plus two
borrowed: (i) a piece of the cut whose top runs \emph{through \(u\)'s parent} is not a
rooted subtree at all, so its size comes from one root-at-\(1\) pass read in Ch.~15's
complement view \(n-\size(u)\) (\(\defref{c20-parent-piece}\)); (ii) the statement's
``identical'' is an edge-preserving bijection with no role for the attachment vertex, so the
key is the unrooted centre form \(\Phi\) (\(\thmref{c22-centreiso}\)), whose centre the code
finds with Ch.~12's two sweeps (\(\thmref{c09-twosweep}\)); (iii) every piece of every
candidate -- and both rootings of a centre edge -- must be interned in \textbf{one} shared
\(\mathit{tid}\) table, or the integers come from different alphabets
(\(\secref{c22-ahu}\), pitfall~3). The reward for filtering first is structural, not
cosmetic: a node whose pieces all have the same size leaves no component bigger than
\(n/2\), i.e.\ it is a \emph{centroid}, and a tree has at most two
(\(\thmref{c09-centexist}\)) -- so the exact stage runs at most twice. Both official
samples are replayed by hand below: the 13-node tree \(\to\) \(3\), the 6-node tree
\(\to\) \(-1\).
\end{notebox}

\solvedproblem{c22-summit}{Regular Forestation}{CF 1252F}
  {https://codeforces.com/problemset/problem/1252/F}{CF 2400 L5}
  {one root-at-\(1\) pass gives every piece size in \(O(1)\) -- child piece \(\size(v)\),
   parent piece \(n-\size(u)\) (\(\defref{c20-parent-piece}\)); the divisibility filter
   \(\deg(u)\mid(n-1)\) plus ``all pieces equal'' admits only centroids, hence at most two
   candidates (\(\thmref{c09-centexist}\)); each piece is then canonicalised \emph{unrooted}
   by two BFS sweeps for the diameter, the centre or centre edge from Ch.~12
   (\(\thmref{c09-twosweep}\)), and the interned AHU labeller (\(\coderef{c22-ahu}\)) run
   against ONE shared \(\mathit{tid}\) table for the whole program, giving
   \(\Phi=(\mathrm{id},\mathrm{id})\) or the sorted pair (\(\thmref{c22-centreiso}\))}

\begin{problemstatement}[Regular Forestation (statement, constraints and both samples verified against the live CF page)]
\(U\) is a tree on \(N\) nodes (\(3\le N\le 4000\)) given by \(N-1\) edges
\((a_i,b_i)\), \(1\le a_i<b_i\le N\). Let \(u\) be a node of degree at least \(2\); removing
\(u\) leaves two or more disconnected smaller trees. \(S\) and \(T\) are called
\textbf{identical} if there is a bijection \(f:V(S)\to V(T)\) such that for all pairs
\((s_i,s_j)\), \(s_i\)--\(s_j\) is an edge of \(S\) \emph{iff} \(f(s_i)\)--\(f(s_j)\) is an
edge of \(T\) -- that is, isomorphic as \emph{unrooted} trees, with no role for how the
piece was attached. A node \(u\) is a \textbf{good cutting point} iff the removal of \(u\)
leaves two or more trees that are \emph{pairwise identical}.

\textbf{Task.} Output the maximum number of disconnected trees obtainable by removing
exactly one good cutting point, or \(-1\) if no good cutting point exists.

\textbf{Input.} One line with \(N\); then \(N-1\) lines, each with \(a_i\ \ b_i\).
\textbf{Output.} One integer. Time limit 1 second, memory 256 MB.

\textbf{Samples (both official).} \#1: \(N{=}13\) with edges
\(1\!-\!5,1\!-\!7,2\!-\!4,2\!-\!8,2\!-\!11,3\!-\!12,4\!-\!7,4\!-\!12,6\!-\!11,7\!-\!13,
9\!-\!10,9\!-\!12\) \(\to\) \texttt{3}. \#2: \(N{=}6\) with edges
\(1\!-\!2,1\!-\!3,2\!-\!4,3\!-\!5,3\!-\!6\) \(\to\) \texttt{-1}.
(The statement itself points out the answer to \#1: node \(4\), giving the three line graphs
\(\{5,1,7,13\}\), \(\{8,2,11,6\}\), \(\{3,12,9,10\}\).)
\end{problemstatement}

\begin{attemptgate}[Attempt it yourself first -- \lv{5}, budget a real sitting]
Two questions before any code: what may you \emph{not} assume about the shape of a piece
(its top vertex is named by the cut, but the definition is not rooted), and what does
``all \(d\) pieces identical'' force on \(d\) and \(N\) immediately?
\peek{If the \(d\) pieces are pairwise identical then they have the same size, so
\(d\cdot s=N-1\): only nodes whose degree divides \(N-1\) can ever be good, and each piece
has size \(s=(N-1)/d\) in \(O(1)\) from one subtree-size pass.}
\nudge{That filter is brutal: a surviving node leaves no component bigger than
\((N-1)/2<N/2\), so it is a centroid, and Ch.~12 proves a tree has at most two
(\(\thmref{c09-centexist}\)). You will therefore have to canonicalise at most two full
cuts -- which makes the exact, collision-free route affordable.}
\fullhint{For a surviving \(u\), each piece is an \emph{unrooted} tree: compute
\(\Phi(\text{piece})\) exactly as \(\coderef{c22-unrooted}\) does -- two BFS sweeps inside
the piece (walk through \(u\) forbidden) for the diameter, take the middle vertex
(\(\Phi\) = that one id, doubled) or the middle \emph{edge} (\(\Phi\) = the sorted pair of
the two endpoint-rooted ids), and label with the interned AHU routine of
\(\coderef{c22-ahu}\). Keep ONE \(\mathit{tid}\) map for the whole program, so ids from
piece \(i\) are comparable with ids from piece \(j\) and from the second rooting of a centre
edge. Answer \(=\max\deg(u)\) over pieces-all-equal \(u\), else \(-1\). Sizes: piece under a
child \(v\) is \(\size(v)\); the piece through \(u\)'s parent is \(N-\size(u)\).}
\end{attemptgate}

\solstep{Understanding}
The input names no root and no child order, and the definition of \emph{identical} is an
edge-preserving bijection, so this is the chapter's cleanest possible shape problem -- with
one twist that has to be handled honestly: the objects being compared are the components of
\(U\setminus\{u\}\), and one of them is the \emph{complement} of a subtree. Two things
follow. First, a candidate is not free: \(d\) pairwise-identical pieces all have size
\(s=(N-1)/d\), so \(\deg(u)\) must divide \(N-1\) and every incident piece must have exactly
that size -- an \(O(\deg)\) test per node, using the \(\size\) table from one DFS and
Ch.~15's complement reading \(N-\size(u)\) for the parent-side piece. Second, no candidate
is arbitrary either: equal pieces of size \(s\le(N-1)/2\) make \(u\) a centroid, so at most
two nodes ever reach the expensive stage. What remains is the exact question ``are these
\(d\) unrooted trees all the same shape?'', answered per piece by \(\Phi\) -- centre from
the diameter (Ch.~12), one shared interned labelling table (Ch.~28.1, 28.3), sorted pair for
a centre edge.

\solstep{Key Observations}
\begin{observation}
\textbf{Equal size is a filter, never a proof -- but it is a very strong filter.}
Identical \(\Rightarrow\) equal size gives \(d\mid(N-1)\) and piece sizes \(s=(N-1)/d\),
readable in \(O(1)\) each from one root-at-\(1\) pass (child piece \(\size(v)\), parent piece
\(N-\size(u)\), \(\defref{c20-parent-piece}\)). Moreover any surviving \(u\) has every
component of size \(s\le(N-1)/2<N/2\), i.e.\ \(u\) is a centroid, and a tree has at most two
centroids, one or two adjacent ones (\(\thmref{c09-centexist}\)). So the number of nodes
whose cut must actually be canonicalised is \(\le2\), and the whole exact stage costs at
most two \(O(N)\) walks. CF~1252F's own \(\le 4000\) budget is generous \emph{because} of
this, not in spite of it.
\end{observation}
\begin{observation}
\textbf{The pieces are unrooted, so the attachment vertex may not enter the key.}
``Identical'' quantifies over all bijections; the node adjacent to \(u\) is distinguished
only by the cut, not by the piece. The certificate is therefore \(\Phi(\text{piece})\) --
the centre-rooted id, or the sorted pair when the piece's centre is an edge
(\(\defref{c22-canroot}\), \(\thmref{c22-centreiso}\)) -- computed with all rootings and all
pieces interned in \textbf{one} \(\mathit{tid}\) table (\(\secref{c22-ahu}\)). Two
consequences to hold on to: the deterministic route is affordable here, so no collision
budget is spent at all (\(\thmref{c22-collision}\) is a tool you choose not to need); and
the same \(\Phi\) code serves every piece, so this problem reuses
\(\coderef{c22-unrooted}\) rather than inventing machinery.
\end{observation}

\solstep{Think It With Me}
\begin{thinkbox}
\deadend{Root each piece at the vertex that touched \(u\) and compare those ids}
{The rule is well defined but answers a \emph{different} question -- rooted isomorphism with
the attachment as root. Take \(N=9\) with edges
\((1,2),(2,3),(3,4),(4,5),(1,7),(6,7),(7,8),(8,9)\): cutting at \(1\) leaves \(4+4\) nodes,
the filter passes (\(d=2\mid 8\), both pieces size \(4\)), and both pieces are the path
\(P_4\). Rooted at the attachment they are different animals: piece \(\{2,3,4,5\}\) is a
chain rooted at an \emph{endpoint} -- interned \(()\!\to\!1,(1)\!\to\!2,(2)\!\to\!3,(3)\!\to\!4\)
gives id \(4\) -- while piece \(\{6,7,8,9\}\) is rooted at a \emph{middle} vertex, giving
\((1,2)\!\to\!3\). \(4\neq3\), so the rule prints \(-1\); the truth is \(2\), because
\(\Phi\) of both pieces is the same sorted pair \((3,3)\). \textbf{False NO -- WA}.}
\deadend{Compare pieces by size only (the filter as if it were the proof)}
{Same shape requires more than equal cardinality. Take \(N=9\) with edges
\((1,2),(2,3),(3,4),(4,5),(1,6),(6,7),(6,8),(6,9)\): cutting at \(1\) leaves the path
\(P_4\) on \(\{2,3,4,5\}\) and the star \(K_{1,3}\) on \(\{6,7,8,9\}\) -- both size \(4\),
both survive the filter. The exact forms disagree: \(P_4\) has a centre edge and
\(\Phi=(3,3)\) (tuple \((1,2)\)), the star has one centre \(6\) and
\(\Phi=(4,4)\) (tuple \((1,1,1)\)). No other node survives the filter either -- at \(6\) the
parent-side piece has size \(N-\size(6)=5\neq 2\), at \(3\) the child piece \(\size(4)=2\neq
4\). Truth \(-1\); the size-only rule prints \(2\). \textbf{False YES -- WA}, and the worst
kind, because it reports a cutting point that does not exist.}
\deadend{One rerooting pass over the commutative shape hash, then read every piece in
\(O(1)\)}
{Two independent failures. (a) The piece through \(u\)'s parent is a \emph{connected
complement}, and a sum/XOR of per-child random values has no subtraction that produces it:
Ch.~15 transfers a fold because the aggregate is invertible or because children are combined
by exclusive prefixes and suffixes (\(\defref{c20-parent-piece}\)) -- not because hashes
reverse. (b) For the pieces that \emph{are} subtrees, \(h(u)\) is the
\(\defref{c22-shapehash}\) value rooted at the attachment -- dead end~1's wrong object with
a probabilistic key stapled on, plus the collision premium of \(\thmref{c22-collision}\) on
a question the judge answers decisively. \textbf{WA}, and the \(O(1)\) promise buys nothing:
after the filter there are at most two candidates to afford the exact route for.}
\aha{Split the work by cost, exactly as the filter suggests. Phase~1 (cheap, exact):
one \(\mathit{buildsz}(1,0)\) pass; then per node, skip unless \(\deg(u)\ge2\),
\(\deg(u)\mid(N-1)\) and every incident piece has size \((N-1)/\deg(u)\) -- at most two
nodes pass, because a passing node is a centroid (\(\thmref{c09-centexist}\)). Phase~2
(expensive, at most twice, exact): for each surviving \(u\) and each neighbour \(v\), walk
inside that piece with \(v\) as the entry and two BFS sweeps for the diameter
(\(\thmref{c09-twosweep}\)), take the centre or centre edge, and intern the rooted form(s)
with the AHU labeller against \textbf{one} global \(\mathit{tid}\) table; \(u\) is good iff
all \(\Phi\) values are equal as sorted pairs. Answer \(=\max\deg(u)\) over good \(u\),
else \(-1\).}
\end{thinkbox}

\solstep{The Approach Ladder}
\begin{center}\renewcommand{\arraystretch}{1.25}\setlength{\tabcolsep}{3pt}
\begin{tabular}{@{}c p{0.24\linewidth} p{0.28\linewidth} l c@{}}
\toprule
\textbf{\#} & \textbf{Approach} & \textbf{Idea} & \textbf{Time} & \textbf{Verdict}\\
\midrule
1 & \lv{2}\ pieces by size, every node & the filter used as a proof & \(\bigO{n^2}\) & \textcolor{warningred}{\ding{55}\ \textbf{WA}} \\
2 & \lv{3}\ attachment-rooted ids & right table, wrong object & \(\bigO{n^2\log n}\) & \textcolor{warningred}{\ding{55}\ \textbf{WA}} \\
3 & \lv{4}\ \(\Phi\) per piece, table rebuilt per piece, no filter & right object, ids from different alphabets & \(\bigO{n^2\log n}\) & \textcolor{warningred}{\ding{55}\ \textbf{WA}} \\
4 & \lv{5}\ \textbf{filter, then \(\Phi\) in ONE shared table} & size as filter, centre form as proof & \(\bigO{n\log n}\) & \ac \\
\bottomrule
\end{tabular}\end{center}
Rungs~1--3 each keep one half of the truth: 1 has the cheap invariant and no certificate, 2
has a certificate for the wrong equivalence, 3 has the right equivalence and breaks the
interning discipline that makes ids meaningful (pitfall~3 in \(\secref{c22-ahu}\)). Rung~4 is
the chapter's thesis assembled: canonicalise once, compare integers, and let a structural
fact (the centroid bound) decide how often you may afford to do it exactly.

\solstep{Solution}
Read the tree. One DFS from node \(1\) fills \(\size\) and \(\parent\) for the whole run.
Scan \(u=1..N\): let \(d=\deg(u)\); require \(d\ge2\), \((N-1)\bmod d=0\), \(s=(N-1)/d\), and
for every neighbour \(v\) the piece size \(\size(v)\) if \(\parent[v]=u\) and \(N-\size(u)\)
otherwise. If any check fails, continue. Otherwise set \(\mathit{blocked}=u\) -- which turns
every walk into a walk \emph{inside one piece} -- and evaluate
\(\Phi=v\!\mapsto\!(\text{two sweeps, centre, labelling})\) for each neighbour, comparing
sorted pairs. The labelling uses the global \(\mathit{tid}\) map, so ids created for
piece~1 are recognised in piece~2 and in the second rooting of a centre edge. Keep
\(\mathit{best}=\max d\) over the \(u\) whose pieces all share one \(\Phi\); print
\(\mathit{best}\) (initialised to \(-1\)).

\solstep{Proof of Correctness}
\begin{proofbox}
\textbf{(Filter is necessary.)} If \(u\) is a good cutting point, its \(d=\deg(u)\ge2\)
pieces are pairwise identical, hence pairwise equal in size, hence \(d\cdot s=N-1\) for the
integer \(s=\size(\text{piece})\); so \(d\mid(N-1)\) and every piece has size \((N-1)/d\),
which the code reads from the frozen \(\size/\parent\) tables (child piece \(=\size(v)\)
because rooting at \(1\) makes \(v\)'s subtree exactly that component; parent piece
\(=N-\size(u)\), Ch.~15's complement). A rejected \(u\) is therefore not good.
\textbf{(Filter is small.)} A surviving \(u\) has every component of size
\(s=(N-1)/d\le(N-1)/2<N/2\), i.e.\ \(u\) is a centroid; a tree has at most two centroids
(\(\thmref{c09-centexist}\)), so at most two nodes reach the exact stage.
\textbf{(Exact stage is right.)} Fix a surviving \(u\). With \(\mathit{blocked}=u\), the
BFS sweeps and \(\mathit{rootify}\) never leave the component of the entry neighbour, so
they compute invariants \emph{of that piece alone}; two sweeps find a diameter
(\(\thmref{c09-twosweep}\)) and its middle is the piece's centre -- one vertex when the
diameter length is even, one edge when odd. The labeller is the interned AHU routine, run on
the piece rooted at that centre (or at each endpoint of the centre edge), so by
\(\thmref{c22-ahu}\) each id is a complete invariant of a rooted shape and, all rootings
sharing one \(\mathit{tid}\) map, ids of different pieces are comparable; by
\(\thmref{c22-centreiso}\), \(\Phi(\text{piece}_i)=\Phi(\text{piece}_j)\) (as a doubled id
or as a sorted pair) \emph{iff} the two pieces are isomorphic as unrooted trees. Hence
``all \(\Phi\) equal'' \(\iff\) ``pairwise identical'' \(\iff\) \(u\) is a good cutting
point, and the recorded count \(d=\deg(u)\) is exactly the number of pieces.
\textbf{(Answer.)} \(\mathit{best}\) is the maximum \(d\) over precisely the good cutting
points, and remains \(-1\) when that set is empty -- the required output. \(\Box\)
\end{proofbox}

\solstep{Implementation}
\begin{lstlisting}[caption={CF 1252F Summit: buildsz once; divisibility plus equal piece sizes keeps at most two candidates (centroids); each piece is canonicalised unrooted by two BFS sweeps, the centre or centre edge, and the interned AHU labeller against ONE global tid table. ASCII only, identifiers without underscores.}, label={lst:c22-summit}]
#include <bits/stdc++.h>
using namespace std;
const int MAXN=4005;
int n, blocked=0, nextId=1;
vector<int> adj[MAXN], kids[MAXN];
int sz[MAXN], par[MAXN], par2[MAXN], cl[MAXN];
int distArr[MAXN], prevArr[MAXN];
map<vector<int>,int> tid;              // ONE table for the entire run

void buildsz(int u,int p){             // frozen tables: blocked stays 0
  sz[u]=1; par[u]=p;
  for(int v : adj[u]) if(v!=p){ buildsz(v,u); sz[u]+=sz[v]; }
}
int bfsFar(int start){                 // stays inside the piece: skips blocked
  for(int i=1;i<=n;i++){ distArr[i]=-1; prevArr[i]=0; }
  queue<int> Q; Q.push(start); distArr[start]=0;
  while(!Q.empty()){
    int u=Q.front(); Q.pop();
    for(int v : adj[u]) if(v!=blocked && distArr[v]<0){
      distArr[v]=distArr[u]+1; prevArr[v]=u; Q.push(v);
    }
  }
  int best=start;
  for(int i=1;i<=n;i++) if(distArr[i]>distArr[best]) best=i;
  return best;
}
void rootify(int u,int p){             // kids[] for the current rooting of one piece
  par2[u]=p; kids[u].clear();
  for(int v : adj[u]) if(v!=p && v!=blocked){ kids[u].push_back(v); rootify(v,u); }
}
int label(int u){                      // interned AHU, shared tid
  vector<int> t;
  for(int v : kids[u]) t.push_back(label(v));
  sort(t.begin(), t.end());            // child order is noise: kill it
  map<vector<int>,int>::iterator it=tid.find(t);
  if(it!=tid.end()) return cl[u]=it->second;
  int id=nextId++; tid[t]=id; return cl[u]=id;
}
pair<int,int> phiOf(int entry){        // UNROOTED form of the piece entered at entry
  int A=bfsFar(entry);                 // sweep 1
  int B=bfsFar(A);                     // sweep 2: A..B is a diameter
  int len=distArr[B], m=B;
  for(int i=0;i<len/2;i++) m=prevArr[m];        // middle vertex, or one end of mid edge
  if(len%2==0){ rootify(m,0); int x=label(m); return make_pair(x,x); }
  int c1=m, c2=prevArr[m];                      // centre EDGE: keep both forms
  rootify(c1,0); int x=label(c1);
  rootify(c2,0); int y=label(c2);
  if(y<x) swap(x,y);
  return make_pair(x,y);
}
int main(){
  scanf("%d",&n);
  for(int i=1;i<n;i++){ int a,b; scanf("%d%d",&a,&b);
    adj[a].push_back(b); adj[b].push_back(a); }
  buildsz(1,0);
  int best=-1;
  for(int u=1;u<=n;u++){
    int d=(int)adj[u].size();
    if(d<2 || (n-1)%d!=0) continue;             // sizes must divide
    int s=(n-1)/d, ok=1;
    for(int v : adj[u]){
      int piece=(par[v]==u) ? sz[v] : n-sz[u];  // parent piece = complement view
      if(piece!=s){ ok=0; break; }
    }
    if(!ok) continue;                           // <= 2 nodes ever get past this
    blocked=u;
    pair<int,int> first=phiOf(adj[u][0]);
    for(int j=1;j<d && ok;j++) if(phiOf(adj[u][j])!=first) ok=0;
    blocked=0;
    if(ok && d>best) best=d;
  }
  printf("%d\n",best);
  return 0;
}
\end{lstlisting}

\solstep{Trace on the Official Samples}
\begin{tracebox}[Sample 1: \(N{=}13\) -- the filter keeps exactly one node, and it is the centroid]
\(\mathit{buildsz}(1,0)\) gives \(\size{=}[\,13,4,1,9,1,1,11,1,2,1,2,4,1\,]\) for nodes
\(1..13\) and \(\parent{=}[0,4,12,7,1,11,1,2,12,9,2,4,7]\). Candidates with
\(\deg(u)\ge2\) and \(\deg(u)\mid 12\): \(\{1,2,4,7,9,11,12\}\), each needing piece size
\(s=12/\deg\).
\(u{=}1\,(d2,s6)\): piece via \(5=\size(5)=1\) -- dead.
\(u{=}2\,(d3,s4)\): \(4\) is \(2\)'s \emph{parent}, piece \(=13-\size(2)=9\) -- dead.
\(u{=}7\,(d3,s4)\): piece via \(1=13-\size(7)=2\) -- dead.
\(u{=}9\,(d2,s6)\): via \(10=1\); \(u{=}11\,(d2,s6)\): via \(2=13-2=11\);
\(u{=}12\,(d3,s4)\): via \(3=1\) -- all dead.
\(u{=}4\,(d3,s4)\): via \(2=\size(2)=4\); via \(7\) (its parent) \(=13-\size(4)=13-9=4\);
via \(12=\size(12)=4\) -- \textbf{the only survivor} (and indeed the centroid).
Now \(\mathit{blocked}=4\), pieces entered in order \(2,7,12\).
Piece via \(2=\{2,8,11,6\}\): sweeps give \(A{=}6,B{=}8,\text{len}=3\) (odd), so
\(m=\mathit{prev}[8]=2\), centre edge \((2,11)\). Rooted at \(2\):
\(()\!\to\!1\) for leaves \(8,6\); \((1)\!\to\!2\) for \(11\); \((1,2)\!\to\!3\) for \(2\).
Rooted at \(11\): again \((1,2)=3\) already interned. \(\Phi=(3,3)\).
Piece via \(7=\{7,1,5,13\}\): \(A{=}5,B{=}13,\text{len}=3\), centre edge \((7,1)\); both
rootings re-use \(1,2,3\): \(\Phi=(3,3)\).
Piece via \(12=\{12,3,9,10\}\): \(A{=}10,B{=}3,\text{len}=3\), centre edge \((12,9)\),
\(\Phi=(3,3)\). All three equal -- and \(\mathit{tid}\) never grew past the three tuples
\(\{\,(),(1),(1,2)\,\}\), because all three pieces are the same \(P_4\).
\(\mathit{best}=\max(-1,3)=\boxed{3}\) -- matches \texttt{3}. \(\Box\)
\end{tracebox}

\begin{tracebox}[Sample 2: \(N{=}6\) -- arithmetic alone closes the case]
Degrees: \(\deg(1)=2,\deg(2)=2,\deg(3)=3\), leaves \(4,5,6\) have \(\deg=1\). Here
\(N-1=5\) is \emph{prime}, so the only divisors are \(1\) and \(5\); no node of degree
\(\ge2\) has degree \(5\), and degrees \(2\) and \(3\) fail \((5\bmod d)\neq0\). Zero nodes
reach the exact stage, so \(\mathit{best}\) stays at its initial value: print
\(\boxed{-1}\) -- matches \texttt{-1}. Note what this sample does \emph{not} test: it never
exercises the shape comparison at all, which is precisely why the two nine-node
counterexamples in the thinkbox are the real specification of this problem. \(\Box\)
\end{tracebox}

\solstep{Complexity}
\begin{complexitybox}
\textbf{Time} \(\bigO{n\log n}\) for the filter and the labelling work, \(\bigO{n^2}\) for
the crude \(\bigO n\) reset inside \code{bfsFar} -- see the note below; the honest bound to
carry is ``one DFS, one linear scan, and at most two canonicalisations''. Concretely:
\(\mathit{buildsz}\) is \(\bigO n\); the filter is \(\bigO{\sum_u\deg(u)}=\bigO n\); per
surviving candidate, each of its \(d\) pieces costs two BFS sweeps \(\bigO m\) plus
\(\mathit{rootify}+\mathit{label}\) at \(\bigO{m\log m}\) for the sorted child tuples and the
\(\mathit{tid}\) lookups, so \(\bigO{n\log n}\) per candidate and \(\le 2\) candidates.
\textbf{Space} \(\bigO n\) for the arrays and the pieces, plus the shared \(\mathit{tid}\)
table, which never holds more tuples than there are distinct rooted shapes interned
(\(\le 2(n-1)\) over a run, since each piece's labelling creates at most one id per vertex
and every id after the first comes from a real shape). At \(N\le4000\) this is a few
megabytes and comfortably inside 1\,s.
\end{complexitybox}

\begin{hardeningbox}
\begin{itemize}[nosep,leftmargin=1.4em]
  \item \textbf{Keep the filter, not just the shapes.} The \(\bigO n\) reset loop in
        \code{bfsFar} makes a degree-\(d\) candidate cost \(\bigO{d\cdot n}\); only the
        centroid bound (\(\le2\) survivors) is what keeps the total near-linear in practice.
        Porting to \(N=2\cdot10^5\)? Replace the reset with a visit token
        \code{seen[i]=stamp}, and the bound then gives \(\bigO{n\log n}\) outright.
  \item \textbf{Freeze \(\mathit{sz}/\mathit{par}\) before you start rooting.} The filter
        reads the tables from the single \(\mathit{buildsz}(1,0)\) pass, so
        \(\mathit{rootify}\) writes into \emph{separate} arrays
        (\(\mathit{par2},\mathit{kids}\)). Sharing one \(\mathit{par}\) between the two
        jobs silently corrupts the piece sizes -- and it corrupts them only for the
        candidates you already like, which is how this bug reaches the judge.
  \item \textbf{Compare the sorted pair, not a packed integer.} \code{min(x,y)*K +
        max(x,y)} overflows for \(K\) large once ids reach \(\sim\!2n\); the sorted
        \code{pair<int,int>} is the honest key, exactly as in
        \(\coderef{c22-unrooted}\).
  \item \textbf{The survivor is a centroid, not ``the centre''.} In sample~1 the good node
        \(4\) also happens to be the tree's diameter centre, and it is tempting to promote
        that to a candidate generator. Do not: the centre is defined by \emph{distances} and
        the centroid by \emph{sizes}, and the equal-size condition is a size statement.
        Generate candidates from the filter, and use the centre only \emph{inside} a piece,
        where \(\thmref{c22-centreiso}\) says it belongs.
  \item \textbf{Stay deterministic.} \(\defref{c22-shapehash}\) would let you compare pieces
        in \(O(1)\) after one \(\bigO n\) pass, but with at most two candidates the exact
        route is already cheap, so there is no \(\le\depth/2^{64}\) risk to accept here --
        the ledger of \(\thmref{c22-collision}\) is spent only when you choose to spend it.
\end{itemize}
\end{hardeningbox}

\begin{recallbox}[Chapter 28 in ten seconds]
\emph{A shape is what survives deleting every accidental freedom: names, child order and
root choice. Intern each vertex's \textbf{sorted} child-id tuple into a fresh id and rooted
isomorphism becomes integer equality (\(\thmref{c22-ahu}\), \(\coderef{c22-ahu}\)); an
unrooted tree roots itself at the centre Ch.~12 hands you, giving \(\Phi\) as one id or as a
\textbf{sorted pair} for a centre edge (\(\thmref{c22-centreiso}\),
\(\coderef{c22-unrooted}\)); a mirror-pairing order exists iff every vertex has at most one
child class of odd multiplicity (\(\thmref{c22-sym}\), CF~1800G); one labelling is also the
census -- \(|\mathit{tid}|\) shapes, \(\mathit{freq}[\mathrm{cl}(v)]\) twins
(\(\thmref{c22-classes}\)); hashes buy \(O(1)\) compares at the price of a ledger that makes
NO certain and YES budgeted (\(\thmref{c22-collision}\)), and CF~601D's distinct downward
strings are counted with the \emph{positional} \(B^{\depth}\) key merged small-to-large
(\(\probref{c22-organic}\)). The Summit CF~1252F (\(\probref{c22-summit}\)) filters a cut by
sizes -- which is a centroid statement -- and proves it with \(\Phi\), once, in one table.}
\selftest{
  \item Why must both trees (and both rootings of a centre edge) share one \(\mathit{tid}\)
        table? \hfill\lv{3}
  \item State \(\Phi\) for a tree of diameter length \(5\), and for one of length \(4\). \hfill\lv{4}
  \item CF~1800G's test: why is checking the root alone not enough, and what is the pairing
        key? \hfill\lv{4}
  \item Which of the two hashes collapses \(\langle A,B\rangle\) into \(\langle B,A\rangle\),
        and which one does CF~601D need? \hfill\lv{5}
  \item In CF~1252F, why does the equal-size filter leave at most two candidates? \hfill\lv{5}
}
\end{recallbox}

\chaptersection{Tool Kit Inventory}

The machine assembled through Chapter~28. Earlier chapters' rows are carried forward;
this chapter's are appended. Pointers only.

\begin{center}
\small
\renewcommand{\arraystretch}{1.25}
\setlength{\tabcolsep}{3pt}
\begin{tabular}{@{}p{0.24\linewidth} l c p{0.22\linewidth} >{\raggedright\arraybackslash\scriptsize}p{0.17\linewidth}@{}}
\toprule
\textbf{Tool} & \textbf{Since} & \textbf{Lv} & \textbf{Trigger phrase} & \textbf{Where} \\
\midrule
\multicolumn{5}{l}{\itshape Carried forward} \\
Centre, diameter by two sweeps & 12 & \lv{3} & ``the middle the tree defines itself'' & \chref{c09}, \thmref{c09-twosweep} \\
At most two centroids & 12 & \lv{4} & ``every piece is small -- how many candidates?'' & \thmref{c09-centexist} \\
Complement size \(n-\size(u)\) & 15 & \lv{4} & ``the part above me, not below'' & \chref{c20}, \defref{c20-parent-piece} \\
Small-to-large set union & 17 & \lv{4} & ``merge per-node sets, rewrite nothing'' & \chref{c28} \\
Bottom-up fold discipline & 13 & \lv{3} & ``children before parent, always'' & \chref{c10} \\
\(\tin/\tout\) containment & 5 & \lv{3} & ``is \(v\) inside my subtree'' & \chref{c24} \\
\midrule
\multicolumn{5}{l}{\itshape Chapter 28 -- new this chapter} \\
Interned AHU labeller & 28.1 & \lv{4} & ``same rooted tree? compare two ids'' & \thmref{c22-ahu}, \coderef{c22-ahu} \\
One-shared-table discipline & 28.1 & \lv{3} & ``ids from different tables cannot meet'' & \secref{c22-ahu} \\
Centre form \(\Phi\), sorted pair & 28.3 & \lv{5} & ``equal \emph{unrooted} trees'' & \thmref{c22-centreiso}, \coderef{c22-unrooted} \\
Odd-class mirror test & 28.3 & \lv{4} & ``can the children pair up?'' & \thmref{c22-sym}, \coderef{c22-symmetree} \\
Commutative shape hash & 28.2 & \lv{4} & ``many \(O(1)\) shape compares'' & \defref{c22-shapehash}, \coderef{c22-shapehash} \\
Positional \(B^{\depth}\) string hash & 28.2 & \lv{5} & ``downward strings; shift = one constant'' & \defref{c22-strhash} \\
Collision ledger & 28.2 & \lv{4} & ``how much may I trust this YES?'' & \thmref{c22-collision} \\
Adopt-and-pour string sets & 28.2 & \lv{5} & ``distinct downward strings per vertex'' & \probref{c22-organic}, \coderef{c22-organic} \\
Shape census \(|\mathit{tid}|\), \(\mathit{freq}\) & 28.4 & \lv{4} & ``how many shapes / how many twins'' & \thmref{c22-classes}, \coderef{c22-classes} \\
Four-test signature, pick protocol & 28.5 & \lv{5} & ``is this a shape problem at all?'' & \defref{c22-signature}, \coderef{c22-kit} \\
Filter-then-\(\Phi\) (Summit) & 1252F & \lv{5} & ``all pieces of one cut identical'' & \probref{c22-summit}, \coderef{c22-summit} \\
\bottomrule
\end{tabular}
\end{center}

\chaptersection{Related Topics}

\begin{itemize}[nosep, leftmargin=1.4em]
  \item \chref{c21} (Ch.~27, Mo and Tarjan) --- \textbf{backward, the baton this chapter
        catches}: Ch.~27's boundary forwarded ``which subtree equals which'' as the question
        its window could not answer; \(\secref{c22-ahu}\) opens by collecting it and answers
        every pair at once with one labelling. Mo compares \emph{counts}; here the
        comparison is of \emph{structures}.
  \item \chref{c30} (Ch.~26, Virtual Trees) --- contrast: a virtual tree folds a value over
        one \emph{marked set} per query; this chapter names a \emph{shape} once for all sets.
        Both need \(\bigl\{\tin/\tout\bigr\}\)-free reasoning about unordered child sets --
        virtual trees keep the labels, so the interning never happens there.
  \item \chref{c09} (Ch.~12, Depth, diameter, centroid, centre) --- \textbf{tool, used
        twice and never re-derived}: the centre of a diameter is the canonical root of an
        unrooted tree (\(\thmref{c22-centreiso}\)) and the two sweeps run \emph{inside each
        piece} in the Summit; ``one or two adjacent centroids'' (\(\thmref{c09-centexist}\))
        is what bounds the Summit's candidate list to two.
  \item \chref{c20} (Ch.~15, Rerooting) --- tool and cautionary tale: the parent-side piece
        has size \(n-\size(u)\) in \(O(1)\) (\(\defref{c20-parent-piece}\)); but the
        transfer works because the aggregate is invertible or combined by exclusive
        prefixes -- shape hashes are neither, which is Summit dead end~3.
  \item \chref{c28} (Ch.~17, DSU on tree / small-to-large) --- tool: CF~601D's per-node
        string sets are merged by exactly that discipline (\(\coderef{c22-organic}\)); also
        the boundary's owner of every \emph{value}-combining per-node fold
        (\(\defref{c22-signature}\), test~1).
  \item \chref{c10} (Ch.~13, Tree DP) --- boundary: weights, counts and validity conditions
        that depend on values (mixed CF~1646D, CF~486D, CF~622E) are bottom-up numeric
        folds; this chapter's fold is over \emph{shapes}, one new sentence only.
  \item \chref{c18} (Ch.~24, Centroid Decomposition) --- boundary: path statistics, pair
        distributions and ``split at the middle of a path'' jobs (CF~914/E, CF~1771D) leave
        the shape toolkit; the Summit borrows only centroid \emph{existence}, not the
        decomposition.
  \item \chref{c16} (Ch.~22, HLD) / \chref{c23} (Ch.~25, Link-Cut Trees) --- boundary: any
        interleaving of cut/link edits with shape questions makes re-canonicalising illegal;
        dynamic trees carry the ordered, labelled structure this chapter deletes.
  \item \chref{c04} (Ch.~4, Representations) --- where ``rooted vs.\ unrooted, ordered vs.
        unordered children'' was first named; \(\defref{c22-iso}\) is that vocabulary made
        into a decision procedure.
  \item \chref{c31} (Ch.~29, Trees Born from Graphs) --- \textbf{next chapter}: the shapes
        here are trees; the next chapter earns them from general graphs -- spanning-tree
        structure, and what changes when the input is no longer known to be acyclic (the
        caveat \(\defref{c22-pick}\) ends this chapter with).
\end{itemize}

\begin{lastword}[Last Word --- Chapter 28]
\emph{``A shape is what is left when you stop caring about everything accidental -- names,
child order, the root you chose. Delete the freedoms one at a time and the last one is the
trick: let the tree choose its own root, its centre, and when the centre is an edge keep
both halves as a sorted pair. Then isomorphism is integer equality, computed once in
\(\bigO{n\log n}\) with no randomness at all; hashing only changes how often you pay, and
buys the collision ledger with it.''}

\medskip
Chapter~27 ended by admitting it could not say whether two subtrees were \emph{the same},
only how many colours they held. Section 28.1 answered with one idea -- sort each vertex's
child ids and intern the tuple -- and the running example paid for it twice: \(R\) labels to
\(4\), and its unrooted twin \(R'\) rooted at vertex~\(1\) labels to \(7\), the wrong-root
phenomenon in one pair of integers (\(\thmref{c22-ahu}\), \(\figref{c22-twin}\)). 28.2
traded determinism for speed and wrote the bill honestly: a commutative fold for shapes, a
depth-weighted positional fold for strings, and a ledger in which NO is certain and YES
costs at most \(\depth/2^{64}\) (\(\thmref{c22-collision}\)) -- spent on CF~601D's distinct
downward strings, merged small-to-large so nothing is ever rewritten
(\(\probref{c22-organic}\)). 28.3 supplied the missing canonical root -- the centre Ch.~12
gives you -- and the sorted pair that rescues the centre-edge case, then read CF~1800G's
mirror question as one odd child class per vertex (\(\thmref{c22-centreiso}\),
\(\thmref{c22-sym}\)). 28.4 noticed that the same table is already a census, and that the
census is root-dependent (\(\thmref{c22-classes}\)). 28.5 drew the pick -- label-free,
order-free, whole-shape, static -- and handed back everything that failed a test: values to
\(\chref{c10}\) and \(\chref{c28}\), order to traversal strings, distances to \(\chref{c09}\)
and \(\chref{c18}\), edits to \(\chref{c16}\) and \(\chref{c23}\) (\(\defref{c22-pick}\)).
The Summit CF~1252F (\(\probref{c22-summit}\)) put a filter in front of all of it: equal
\emph{sizes} is arithmetic, and only then does \(\Phi\) decide -- with the centroid bound
making the expensive step happen at most twice.

\medskip
Carry three reflexes forward:
\begin{enumerate}[nosep, leftmargin=1.6em]
  \item \emph{Ask what is accidental before what is equal.} Names, child order and the root
        are the three freedoms; each one you can prove irrelevant lets you delete it -- and
        each one that matters belongs to another chapter (\(\defref{c22-signature}\)).
  \item \emph{Canonicalise once, compare forever.} One postorder pass plus one shared table
        answers every pair, counts every class and prices every twin; re-running the
        labeller per query is the \(\bigO n\) tax on a \(\bigO 1\) fact
        (\(\thmref{c22-ahu}\), \(\thmref{c22-classes}\)).
  \item \emph{Unrooted means ``at the centre'', and a centre edge means ``both''.} Any rule
        that names a vertex by label can be broken by renaming; the sorted pair is the
        smallest object that is a function of the tree (\(\thmref{c22-centreiso}\),
        \(\figref{c22-pair}\)).
\end{enumerate}

\medskip
And three habits: keep \emph{one} \(\mathit{tid}\) (or one hash seed) for a whole run and
say so in a comment, because ids from different alphabets compare cleanly and wrongly;
iterate rather than recurse once \(n\) reaches \(3\cdot10^5\) and keep
\(f(v)+c_v\) in \code{long long} (\(\coderef{c22-organic}\)); and when a judge's YES is
decisive, either pay for the exact route or re-verify a hash-YES with an explicit bijection
DFS -- never spend the collision budget twice by accident (\(\thmref{c22-collision}\)). You
now own the equality test that every counting tool before it quietly assumed: given two
trees, you can say in near-linear time whether they are the same tree -- and given any
specification, whether that question is even the right one.

\medskip
\textbf{Next:} \chref{c31} (Chapter 29, Trees Born from Graphs) changes the ground under the
tools. Everything in this chapter assumed the input \emph{is} a tree; the next chapter earns
trees from general graphs -- connectivity, spanning forests, bridges and biconnected
compression -- and asks which structural verdicts survive when a cycle appears. The caveat
\(\defref{c22-pick}\) leaves you with (colour refinement, \(k\)-WL) is the preview: on trees
the canonical form is exact, on graphs it is only the first approximation.
\end{lastword}
